\documentclass[11pt,a4paper]{article}
\usepackage[utf8]{inputenc}
\usepackage[T1]{fontenc}
\usepackage[spanish,es-minimal,es-noshorthands]{babel}
\usepackage{lmodern}
\usepackage[sfdefault,lining]{FiraSans}
\usepackage[a4paper,top=2.3cm,bottom=2.2cm,left=1.8cm,right=1.8cm,headheight=1.3cm,headsep=0.4cm,footskip=1cm]{geometry}
\usepackage{amsmath,amssymb}
\usepackage{icomma}
\usepackage{graphicx}
\usepackage[table]{xcolor}
\usepackage{tikz}
\usepackage{fontawesome5}
\usepackage{ccicons}
\usepackage{fancyhdr}
\usepackage{lastpage}
\usepackage{tasks}
\usepackage{needspace}
\usepackage[most]{tcolorbox}
\usepackage{array}
\usepackage[hidelinks]{hyperref}
\AtBeginDocument{\hypersetup{pdftitle={Ficha de ejercicios de la unidad 1: Números reales. 1º Bach.},pdfsubject={Matemáticas I 1º Bach.},pdfkeywords={números reales, fracción generatriz, valor absoluto, entornos, intervalos, potencias, radicales, racionalización, logaritmos}}}

\definecolor{navy}{HTML}{1B2748}
\definecolor{accent}{HTML}{E8603A}
\definecolor{teal}{HTML}{1F9E90}
\definecolor{soft}{HTML}{EEF1F8}
\definecolor{muted}{HTML}{5D6685}

\setlength{\parindent}{0pt}
\setlength{\parskip}{3pt}
\settasks{label=\alph*), label-format=\bfseries\color{navy}, label-width=1.3em, item-indent=1.8em, column-sep=0.6em, after-item-skip=5pt, before-skip=3pt, after-skip=4pt}

% ---------- capçalera i peu ----------
\newcommand{\llicencia}{\href{https://creativecommons.org/licenses/by-nc-sa/4.0/deed.es}{\ccbyncsa}\ \ Licencia Creative Commons Atribución-NoComercial-CompartirIgual 4.0 Internacional (CC BY-NC-SA 4.0)}
\pagestyle{fancy}
\fancyhf{}
\fancyhead[L]{\raisebox{-0.35\height}{\includegraphics[height=0.95cm]{logo-gva.png}}}
\fancyhead[C]{\small\color{muted}Ficha de ejercicios · Unidad 1: Números reales · 1\textsuperscript{o} Bach.}
\fancyhead[R]{\raisebox{-0.35\height}{\includegraphics[height=0.95cm]{logo-vermellar.png}}}
\fancyfoot[L]{\scriptsize\color{muted}\llicencia}
\fancyfoot[R]{\small\color{muted}\thepage}
\renewcommand{\headrulewidth}{0.4pt}
\renewcommand{\headrule}{\color{muted!40}\hrule height \headrulewidth}
\fancypagestyle{primera}{\fancyhf{}\renewcommand{\headrulewidth}{0pt}
  \fancyfoot[L]{\scriptsize\color{muted}\llicencia}
  \fancyfoot[R]{\small\color{muted}\thepage}}

% ---------- nivells i exercicis ----------
\newcommand{\estrelles}[1]{{\color{accent}\ifcase#1\or\faStar\faStar[regular]\faStar[regular]\or\faStar\faStar\faStar[regular]\or\faStar\faStar\faStar\fi}}
\newcounter{exer}
\newcommand{\ex}[1]{\par\needspace{4\baselineskip}\medskip\refstepcounter{exer}\noindent{\color{navy}\bfseries\theexer.}\ \ }

\newcommand{\seccio}[2]{\par\needspace{16\baselineskip}\bigskip
  \begin{tcolorbox}[colback=navy,colframe=navy,coltext=white,arc=2mm,boxrule=0pt,left=3mm,right=3mm,top=1.2mm,bottom=1.2mm]
  {\large\bfseries #1\quad #2}
  \end{tcolorbox}\vspace{-2pt}}

\newcommand{\exf}[1]{\par\medskip\refstepcounter{exer}\noindent{\color{navy}\bfseries\theexer.}\ \ }
\newcolumntype{C}[1]{>{\centering\arraybackslash}p{#1}}
% quadre de recordatori (només a la fitxa d'exercicis)
\newcommand{\recorda}[1]{\par\smallskip\begin{tcolorbox}[colback=accent!7,colframe=accent!55,arc=2mm,boxrule=0.5pt,left=2.5mm,right=2.5mm,top=1mm,bottom=1mm]\small{\color{accent}\faLightbulb[regular]}\ \textbf{Recuerda:}\ #1\end{tcolorbox}}
\newcommand{\res}[1]{#1}\newcommand{\rest}[1]{#1}

\def\fmneed{11}% que el títol d'una secció no quede sol al final del full
% =====================================================================
% Capçalera FentMates de les fitxes: un color per a cada curs.
% Els colors són els de la pàgina inicial (LEVELS a assets/js/materials.js);
% el logo ve de _eines/logo/pdf/ (python3 _eines/logo/genera.py).
%
% Ús, al final del preàmbul de la unitat:
%   % =====================================================================
% Capçalera FentMates de les fitxes: un color per a cada curs.
% Els colors són els de la pàgina inicial (LEVELS a assets/js/materials.js);
% el logo ve de _eines/logo/pdf/ (python3 _eines/logo/genera.py).
%
% Ús, al final del preàmbul de la unitat:
%   % =====================================================================
% Capçalera FentMates de les fitxes: un color per a cada curs.
% Els colors són els de la pàgina inicial (LEVELS a assets/js/materials.js);
% el logo ve de _eines/logo/pdf/ (python3 _eines/logo/genera.py).
%
% Ús, al final del preàmbul de la unitat:
%   \input{../comu/fentmates.tex}
%   \fmcurs{1eso}
% i, al començament del document:
%   \capcalera{Matemàtiques}{1r ESO}{Fitxa d'exercicis · Unitat 2}{Divisibilitat}
% Opcional abans de \input: \def\fmneed{14} (espai mínim abans de cada secció).
% =====================================================================
\makeatletter
\def\fm@logodir{../../logo/pdf}
\providecommand{\fmneed}{10}
\def\fm@cursnom{}

% clau del curs, color i color del text damunt del color
\newcommand{\fm@curs}[3]{\@namedef{fm@color@#1}{#2}\@namedef{fm@text@#1}{#3}}
\fm@curs{1eso}{E6B84A}{navy}
\fm@curs{2eso}{3B82F6}{white}
\fm@curs{3eso}{FF7A50}{white}
\fm@curs{4eso-a}{A78BFA}{white}
\fm@curs{4eso-b}{6D28D9}{white}
\fm@curs{1batx-mi}{9B1C1C}{white}
\fm@curs{1batx-ccss}{EF6B6B}{white}
\fm@curs{2batx-mii}{1E6B47}{white}
\fm@curs{2batx-ccss}{2BB3A3}{white}

\newcommand{\fmcurs}[1]{%
  \@ifundefined{fm@color@#1}{\PackageError{fentmates}{Curs desconegut: #1}{}}{}%
  \edef\fm@tmp{\noexpand\definecolor{curs}{HTML}{\@nameuse{fm@color@#1}}%
               \noexpand\colorlet{curstext}{\@nameuse{fm@text@#1}}}\fm@tmp
  \colorlet{cursfosc}{curs!72!black}%
  \def\fm@logo{\fm@logodir/vertical-#1.pdf}%
  \def\fm@marca{\fm@logodir/marca-#1.pdf}}

% ---------- capçalera de la primera pàgina ----------
% #1 matèria (petit), #2 curs, #3 tipus de document i unitat, #4 títol
\newcommand{\capcalera}[4]{%
  \gdef\fm@cursnom{#2}%
  \par\noindent
  \begin{tcolorbox}[colback=curs,colframe=curs,boxrule=0pt,arc=2.2mm,width=3.5cm,height=2.15cm,
      valign=center,left=3mm,right=2mm,top=1mm,bottom=1mm,nobeforeafter,box align=center]
    \color{curstext}{\footnotesize #1}\par\vspace{1pt}{\fontsize{19}{21}\selectfont\bfseries #2\par}%
  \end{tcolorbox}\hspace{5mm}%
  \begin{minipage}[c]{\dimexpr\linewidth-3.5cm-5mm-2.9cm\relax}
    {\color{cursfosc}\small\bfseries\MakeUppercase{#3}}\par\vspace{2pt}
    {\color{navy}\fontsize{20}{23}\selectfont\bfseries #4\par}
  \end{minipage}\hfill
  \begin{minipage}[c]{2.6cm}\raggedleft\includegraphics[width=2.6cm]{\fm@logo}\end{minipage}\par
  \vspace{2pt}}

% ---------- barres de secció del color del curs ----------
\renewcommand{\seccio}[2]{\par\needspace{\fmneed\baselineskip}\bigskip
  \begin{tcolorbox}[colback=curs,colframe=curs,coltext=curstext,arc=2mm,boxrule=0pt,left=3mm,right=3mm,top=1.2mm,bottom=1.2mm]
  {\large\bfseries #1\quad #2}
  \end{tcolorbox}\vspace{-2pt}}

% ---------- estrelles de dificultat del color del curs (a Batxillerat no se'n posen) ----------
\@ifundefined{estrelles}{}{%
  \renewcommand{\estrelles}[1]{{\color{curs}\ifcase#1\or\faStar\faStar[regular]\faStar[regular]\or\faStar\faStar\faStar[regular]\or\faStar\faStar\faStar\fi}}}

% ---------- capçalera i peu de les pàgines ----------
% primera pàgina: logos institucionals al peu
\fancypagestyle{primera}{\fancyhf{}\renewcommand{\headrulewidth}{0pt}
  \fancyfoot[L]{\raisebox{-0.32\height}{\includegraphics[height=0.78cm]{logo-gva.png}}\hspace{5mm}%
    \raisebox{-0.32\height}{\includegraphics[height=0.78cm]{logo-vermellar.png}}}
  \fancyfoot[R]{\scriptsize\color{muted}\href{https://creativecommons.org/licenses/by-nc-sa/4.0/deed.ca}{\ccbyncsa}\ \ CC BY-NC-SA 4.0 · fentmates.org\qquad\small\thepage}}
% la resta: el símbol a l'esquerra i el curs a la dreta, amb el títol al centre com abans
\fancyhead[L]{\raisebox{-0.3\height}{\includegraphics[height=0.62cm]{\fm@marca}}}
\fancyhead[R]{\small\bfseries\color{cursfosc}\fm@cursnom}
\renewcommand{\headrule}{\color{curs!45}\hrule height \headrulewidth}
\makeatother

%   \fmcurs{1eso}
% i, al començament del document:
%   \capcalera{Matemàtiques}{1r ESO}{Fitxa d'exercicis · Unitat 2}{Divisibilitat}
% Opcional abans de \input: \def\fmneed{14} (espai mínim abans de cada secció).
% =====================================================================
\makeatletter
\def\fm@logodir{../../logo/pdf}
\providecommand{\fmneed}{10}
\def\fm@cursnom{}

% clau del curs, color i color del text damunt del color
\newcommand{\fm@curs}[3]{\@namedef{fm@color@#1}{#2}\@namedef{fm@text@#1}{#3}}
\fm@curs{1eso}{E6B84A}{navy}
\fm@curs{2eso}{3B82F6}{white}
\fm@curs{3eso}{FF7A50}{white}
\fm@curs{4eso-a}{A78BFA}{white}
\fm@curs{4eso-b}{6D28D9}{white}
\fm@curs{1batx-mi}{9B1C1C}{white}
\fm@curs{1batx-ccss}{EF6B6B}{white}
\fm@curs{2batx-mii}{1E6B47}{white}
\fm@curs{2batx-ccss}{2BB3A3}{white}

\newcommand{\fmcurs}[1]{%
  \@ifundefined{fm@color@#1}{\PackageError{fentmates}{Curs desconegut: #1}{}}{}%
  \edef\fm@tmp{\noexpand\definecolor{curs}{HTML}{\@nameuse{fm@color@#1}}%
               \noexpand\colorlet{curstext}{\@nameuse{fm@text@#1}}}\fm@tmp
  \colorlet{cursfosc}{curs!72!black}%
  \def\fm@logo{\fm@logodir/vertical-#1.pdf}%
  \def\fm@marca{\fm@logodir/marca-#1.pdf}}

% ---------- capçalera de la primera pàgina ----------
% #1 matèria (petit), #2 curs, #3 tipus de document i unitat, #4 títol
\newcommand{\capcalera}[4]{%
  \gdef\fm@cursnom{#2}%
  \par\noindent
  \begin{tcolorbox}[colback=curs,colframe=curs,boxrule=0pt,arc=2.2mm,width=3.5cm,height=2.15cm,
      valign=center,left=3mm,right=2mm,top=1mm,bottom=1mm,nobeforeafter,box align=center]
    \color{curstext}{\footnotesize #1}\par\vspace{1pt}{\fontsize{19}{21}\selectfont\bfseries #2\par}%
  \end{tcolorbox}\hspace{5mm}%
  \begin{minipage}[c]{\dimexpr\linewidth-3.5cm-5mm-2.9cm\relax}
    {\color{cursfosc}\small\bfseries\MakeUppercase{#3}}\par\vspace{2pt}
    {\color{navy}\fontsize{20}{23}\selectfont\bfseries #4\par}
  \end{minipage}\hfill
  \begin{minipage}[c]{2.6cm}\raggedleft\includegraphics[width=2.6cm]{\fm@logo}\end{minipage}\par
  \vspace{2pt}}

% ---------- barres de secció del color del curs ----------
\renewcommand{\seccio}[2]{\par\needspace{\fmneed\baselineskip}\bigskip
  \begin{tcolorbox}[colback=curs,colframe=curs,coltext=curstext,arc=2mm,boxrule=0pt,left=3mm,right=3mm,top=1.2mm,bottom=1.2mm]
  {\large\bfseries #1\quad #2}
  \end{tcolorbox}\vspace{-2pt}}

% ---------- estrelles de dificultat del color del curs (a Batxillerat no se'n posen) ----------
\@ifundefined{estrelles}{}{%
  \renewcommand{\estrelles}[1]{{\color{curs}\ifcase#1\or\faStar\faStar[regular]\faStar[regular]\or\faStar\faStar\faStar[regular]\or\faStar\faStar\faStar\fi}}}

% ---------- capçalera i peu de les pàgines ----------
% primera pàgina: logos institucionals al peu
\fancypagestyle{primera}{\fancyhf{}\renewcommand{\headrulewidth}{0pt}
  \fancyfoot[L]{\raisebox{-0.32\height}{\includegraphics[height=0.78cm]{logo-gva.png}}\hspace{5mm}%
    \raisebox{-0.32\height}{\includegraphics[height=0.78cm]{logo-vermellar.png}}}
  \fancyfoot[R]{\scriptsize\color{muted}\href{https://creativecommons.org/licenses/by-nc-sa/4.0/deed.ca}{\ccbyncsa}\ \ CC BY-NC-SA 4.0 · fentmates.org\qquad\small\thepage}}
% la resta: el símbol a l'esquerra i el curs a la dreta, amb el títol al centre com abans
\fancyhead[L]{\raisebox{-0.3\height}{\includegraphics[height=0.62cm]{\fm@marca}}}
\fancyhead[R]{\small\bfseries\color{cursfosc}\fm@cursnom}
\renewcommand{\headrule}{\color{curs!45}\hrule height \headrulewidth}
\makeatother

%   \fmcurs{1eso}
% i, al començament del document:
%   \capcalera{Matemàtiques}{1r ESO}{Fitxa d'exercicis · Unitat 2}{Divisibilitat}
% Opcional abans de \input: \def\fmneed{14} (espai mínim abans de cada secció).
% =====================================================================
\makeatletter
\def\fm@logodir{../../logo/pdf}
\providecommand{\fmneed}{10}
\def\fm@cursnom{}

% clau del curs, color i color del text damunt del color
\newcommand{\fm@curs}[3]{\@namedef{fm@color@#1}{#2}\@namedef{fm@text@#1}{#3}}
\fm@curs{1eso}{E6B84A}{navy}
\fm@curs{2eso}{3B82F6}{white}
\fm@curs{3eso}{FF7A50}{white}
\fm@curs{4eso-a}{A78BFA}{white}
\fm@curs{4eso-b}{6D28D9}{white}
\fm@curs{1batx-mi}{9B1C1C}{white}
\fm@curs{1batx-ccss}{EF6B6B}{white}
\fm@curs{2batx-mii}{1E6B47}{white}
\fm@curs{2batx-ccss}{2BB3A3}{white}

\newcommand{\fmcurs}[1]{%
  \@ifundefined{fm@color@#1}{\PackageError{fentmates}{Curs desconegut: #1}{}}{}%
  \edef\fm@tmp{\noexpand\definecolor{curs}{HTML}{\@nameuse{fm@color@#1}}%
               \noexpand\colorlet{curstext}{\@nameuse{fm@text@#1}}}\fm@tmp
  \colorlet{cursfosc}{curs!72!black}%
  \def\fm@logo{\fm@logodir/vertical-#1.pdf}%
  \def\fm@marca{\fm@logodir/marca-#1.pdf}}

% ---------- capçalera de la primera pàgina ----------
% #1 matèria (petit), #2 curs, #3 tipus de document i unitat, #4 títol
\newcommand{\capcalera}[4]{%
  \gdef\fm@cursnom{#2}%
  \par\noindent
  \begin{tcolorbox}[colback=curs,colframe=curs,boxrule=0pt,arc=2.2mm,width=3.5cm,height=2.15cm,
      valign=center,left=3mm,right=2mm,top=1mm,bottom=1mm,nobeforeafter,box align=center]
    \color{curstext}{\footnotesize #1}\par\vspace{1pt}{\fontsize{19}{21}\selectfont\bfseries #2\par}%
  \end{tcolorbox}\hspace{5mm}%
  \begin{minipage}[c]{\dimexpr\linewidth-3.5cm-5mm-2.9cm\relax}
    {\color{cursfosc}\small\bfseries\MakeUppercase{#3}}\par\vspace{2pt}
    {\color{navy}\fontsize{20}{23}\selectfont\bfseries #4\par}
  \end{minipage}\hfill
  \begin{minipage}[c]{2.6cm}\raggedleft\includegraphics[width=2.6cm]{\fm@logo}\end{minipage}\par
  \vspace{2pt}}

% ---------- barres de secció del color del curs ----------
\renewcommand{\seccio}[2]{\par\needspace{\fmneed\baselineskip}\bigskip
  \begin{tcolorbox}[colback=curs,colframe=curs,coltext=curstext,arc=2mm,boxrule=0pt,left=3mm,right=3mm,top=1.2mm,bottom=1.2mm]
  {\large\bfseries #1\quad #2}
  \end{tcolorbox}\vspace{-2pt}}

% ---------- estrelles de dificultat del color del curs (a Batxillerat no se'n posen) ----------
\@ifundefined{estrelles}{}{%
  \renewcommand{\estrelles}[1]{{\color{curs}\ifcase#1\or\faStar\faStar[regular]\faStar[regular]\or\faStar\faStar\faStar[regular]\or\faStar\faStar\faStar\fi}}}

% ---------- capçalera i peu de les pàgines ----------
% primera pàgina: logos institucionals al peu
\fancypagestyle{primera}{\fancyhf{}\renewcommand{\headrulewidth}{0pt}
  \fancyfoot[L]{\raisebox{-0.32\height}{\includegraphics[height=0.78cm]{logo-gva.png}}\hspace{5mm}%
    \raisebox{-0.32\height}{\includegraphics[height=0.78cm]{logo-vermellar.png}}}
  \fancyfoot[R]{\scriptsize\color{muted}\href{https://creativecommons.org/licenses/by-nc-sa/4.0/deed.ca}{\ccbyncsa}\ \ CC BY-NC-SA 4.0 · fentmates.org\qquad\small\thepage}}
% la resta: el símbol a l'esquerra i el curs a la dreta, amb el títol al centre com abans
\fancyhead[L]{\raisebox{-0.3\height}{\includegraphics[height=0.62cm]{\fm@marca}}}
\fancyhead[R]{\small\bfseries\color{cursfosc}\fm@cursnom}
\renewcommand{\headrule}{\color{curs!45}\hrule height \headrulewidth}
\makeatother

\fmcurs{1batx-mi}

\begin{document}
\thispagestyle{primera}

\capcalera{Matemáticas I}{1\textsuperscript{o} Bach.}{Ficha de ejercicios · Unidad 1}{\raggedright Números reales}
\vspace{10pt}
Nombre y apellidos: \rule{7.8cm}{0.4pt}\hfill Grupo: \rule{1.6cm}{0.4pt}\hfill Fecha: \rule{2.2cm}{0.4pt}
\vspace{8pt}

\begin{tcolorbox}[colback=soft,colframe=soft,arc=2mm,left=3mm,right=3mm,top=1.5mm,bottom=1.5mm]
\small Resuelve los ejercicios en el cuaderno mostrando todos los pasos. Da los resultados de manera exacta (fracciones y raíces simplificadas, sin raíces en el denominador) y, si hace falta, también aproximados. Los ejercicios 1--42 son los de la ficha del departamento; la sección 6 tiene otros nuevos para ampliar y aplicar lo que has aprendido.
\end{tcolorbox}

% =====================================================
\seccio{1}{Clasificación de los números}

\exf{1} Indica el menor de los conjuntos numéricos a los que pertenecen los números siguientes:\begin{tasks}(5)
  \task $\sqrt5$
  \task $2$
  \task $-3$
  \task $24/6$
  \task $2,45$
  \task $-1,33$
  \task $2,3\widehat{45}$
  \task $\pi$
  \task $5/12$
  \task $0$
\end{tasks}

\ex{1} Razona si las afirmaciones siguientes son verdaderas o falsas:\begin{tasks}(1)
  \task Todo número entero es racional.
  \task Todo número decimal se puede expresar en forma de fracción.
  \task Todo número natural se puede escribir como una fracción.
  \task Si se suman dos números racionales, el resultado es siempre un número racional.
  \task Si se multiplican dos números racionales no enteros, el resultado es siempre un número racional no entero.
\end{tasks}

\ex{2} Realiza las operaciones siguientes escribiendo previamente los números en forma de fracción:\begin{tasks}(2)
  \task $2,\widehat7+3,\widehat4-1,\widehat1$
  \task $1,\widehat8+0,\widehat1$
  \task $1,2\widehat3+4,572-2,5\widehat{67}$
  \task $3,\widehat{45}\cdot1,23\widehat4-2,5$
\end{tasks}\recorda{Fracción generatriz: en el numerador, el número sin coma menos la parte no periódica; en el denominador, tantos 9 como cifras tiene el período y tantos 0 como cifras tiene el anteperíodo.}


% =====================================================
\seccio{2}{Valor absoluto}

\exf{2} Averigua para qué valores de $x$ se cumplen las desigualdades siguientes:\begin{tasks}(3)
  \task $|x|<3$
  \task $|x|\ge3$
  \task $|x-2|\le3$
  \task $|x|\le5$
  \task $|x-4|<2$
  \task $|x-4|\le2$
  \task $|x-4|>2$
  \task $|2x-1|>1$
  \task $|-5x+3|<7$
  \task $\left|\frac{3x}{2}-6\right|>10$
  \task $\left|1-\frac{2x}{5}\right|\le3$
  \task $|-5x+25|\le-2$
  \task $\left|x+\frac34\right|>-1$
\end{tasks}\recorda{Si $r>0$: \ $|A|<r\iff -r<A<r$ \quad y \quad $|A|>r\iff A<-r$ o $A>r$. Al multiplicar o dividir por un número negativo, la desigualdad cambia de sentido.}


% =====================================================
\seccio{3}{Intervalos y entornos}

\exf{1} Completa la tabla siguiente:

\begin{center}\renewcommand{\arraystretch}{1.5}\begin{tabular}{|C{4.6cm}|C{3.2cm}|C{6cm}|}\hline\rowcolor{soft}\textbf{Desigualdad} & \textbf{Intervalo} & \textbf{Representación}\\\hline$\{x\in\mathbb R\ /\ x<-3\}$ & & \\\hline & $x\in[-2,\, 4)$ & \\\hline $\{x\in\mathbb R\ /\ -5<x\le-3\}$ & & \\\hline & $x\in(4,\, \infty)$ & \\\hline\end{tabular}\end{center}

\ex{1} Representa los conjuntos numéricos siguientes:\begin{tasks}(3)
  \task Números mayores que 3.
  \task $\{x\in\mathbb R/2\le x\le5\}$
  \task $\{x\in\mathbb R/3<x\le7\}$
  \task Números menores que 1 excluyendo el 0.
  \task $(-3,\, -1)$
  \task $[4,\, +\infty)$
  \task Números reales positivos.
  \task $\{x\in\mathbb R/-2\le x<5\}$
  \task $3x-2\le7$
  \task $-4x+1>8$
  \task $|x|<5$
  \task $|5x-6|\ge3$
\end{tasks}

\ex{1} Expresa como desigualdad, como intervalo y representa:\begin{tasks}(2)
  \task $x$ es menor que $-5$.
  \task 3 es menor o igual que $x$.
  \task $x$ está comprendido entre $-5$ y 1.
  \task $x$ está entre $-2$ y 0, ambos incluidos.
\end{tasks}

\ex{1} Escribe la desigualdad que verifica todo número $x$ que pertenece a estos intervalos y represéntalos gráficamente:\begin{tasks}(4)
  \task $[-2,\, 7]$
  \task $[13,\, +\infty)$
  \task $(-\infty,\, 0)$
  \task $(-3,\, 0]$
  \task $\left[\frac32,\, 6\right)$
  \task $(-\infty,\, +\infty)$
  \task $(1,\, 6)$
  \task $(0,\, 4]$
  \task $(-\infty,\, 2)$
  \task $\left[2,\, \frac43\right]$
  \task $(-1,\, +\infty)$
  \task $\left(2,\, \frac72\right)$
\end{tasks}

\ex{2} Dados los intervalos $A=[-3,\, 2]$, $B=[0,\, 5]$, $C=[2,\, +\infty)$, $D=(0,\, 10)$, $E=\left(2,\, \frac72\right)$ y $F=\left(1,\, \frac52\right)$, calcula:\begin{tasks}(3)
  \task $A\cap B$
  \task $A\cap D$
  \task $A\cup B$
  \task $(A\cup B)\cap C$
  \task $E\cap F$
  \task $F-C$
  \task $(B-E)\cap A$
  \task $(A\cup B)\cap(E\cap F)$
  \task $(A\cap B)-(D\cup F)$
\end{tasks}

\ex{2} Dados los intervalos $A=\,]-\infty,\, -4]$, $B=[2,\, 5[$, $C=[-6,\, 1]$, $D=[0,\, +\infty[$ y $E=\,]-2,\, -1[$, calcula:\begin{tasks}(3)
  \task $A\cup B$
  \task $B\cap C$
  \task $D\cap C$
  \task $A\cup(B\cap D)$
  \task $B\cap(A\cup E)$
  \task $B\cap D\cap E$
  \task $(B\cup C)\cap A$
  \task $D-B$
  \task $C-E$
  \task $(A\cup C)-(C\cap E)$
\end{tasks}

\ex{1} Escribe en forma de intervalo los números que verifican las desigualdades siguientes:\begin{tasks}(4)
  \task $x<3$ y $x\ge5$
  \task $x>0$ y $x<4$
  \task $x\le-1$ y $x>1$
  \task $x<3$ y $x\ge-2$
\end{tasks}


% =====================================================
\seccio{4}{Potencias y raíces}

\exf{1} Calcula las potencias siguientes:\begin{tasks}(5)
  \task $2^3$
  \task $(-3)^2$
  \task $-3^2$
  \task $(-0,15)^2$
  \task $(0,01)^4$
  \task $\left(\frac15\right)^2$
  \task $\left(-\frac32\right)^2$
  \task $\left(-\frac23\right)^3$
  \task $\left(-\frac13\right)^4$
  \task $(-0,003)^5$
\end{tasks}

\ex{1} Reduce las expresiones a una sola potencia, aplicando las propiedades de las potencias:\begin{tasks}(2)
  \task $(-3)^2\cdot(-3)^3\cdot(-3)$
  \task $\left(\frac12\right)^2\cdot\left(\frac12\right)^5\cdot\left(\frac12\right)^3\cdot\frac12$
  \task $[(-3)^2]^3$
  \task $(-3^2)^3$
  \task $(-1^4)^3$
  \task $\left[\left(\frac12\right)^2\right]^5$
  \task $\{[(-3^2)^3]^5\}^3$
  \task $\{[(-0,1)^2]^3\}^3$
  \task $(-3)^6:(-3)^6$
  \task $(-4)^5:(-4)^3$
  \task $\left(-\frac23\right)^3:\left(-\frac23\right)$
  \task $(\sqrt x)^5:(\sqrt x)^3$
  \task $\left(\frac13\right)^{-2}\cdot3^7\cdot3^{-8}\cdot\left(\frac13\right)^3$
  \task $\left[\left(\frac23\right)^5\right]^{-1}\cdot\left(\frac32\right)\cdot\left(\frac23\right)^{-3}$
\end{tasks}

\ex{2} Simplifica al máximo las expresiones siguientes:\begin{tasks}(3)
  \task $\dfrac{2^7\cdot2^{-8}}{2^4\cdot2^5}$
  \task $\dfrac{a^4\cdot(b\cdot c)^6\cdot a^{-2}}{b^3\cdot c^5}$
  \task $\dfrac{(p\cdot q)^3\cdot(p^2\cdot q^3)^2}{(p\cdot q)^{-2}}$
  \task $\dfrac{32^6\cdot27^4\cdot30^3}{10^5\cdot15^{14}}$
  \task $\dfrac18\cdot48^3\cdot9^2\cdot36^5$
  \task $\dfrac{2^{-3}\cdot5^{\frac12}\cdot(2^4)^{\frac12}}{5^{-\frac32}\cdot2^{-5}}$
\end{tasks}

\ex{2} Calcula de dos formas (operando primero dentro del corchete, o aplicando las propiedades) las potencias siguientes:\begin{tasks}(3)
  \task $\left[\left(\frac13\right)^2\cdot\left(\frac35\right)\right]^3$
  \task $\left(\frac34:\frac56\right)^3$
  \task $\left[\dfrac{3^2\cdot(-2)}{6}\right]^2$
  \task $\left[\left(-\frac13\right)^3\cdot\left(\frac65\right)\cdot(-5)\right]^4$
  \task $\left[\left(-\frac32\right):\left(\frac34\right)\right]^5$
  \task $\left[\dfrac{(-2)\cdot(-3)^3}{9}\right]^2$
\end{tasks}

\ex{2} Calcula el valor de las expresiones siguientes:\begin{tasks}(2)
  \task $\left(\frac34\right)^2-\frac58$
  \task $\left(-\frac12\right)^2+\left(\frac32\right)^3-\left(\frac53\right)^2$
  \task $\left(\frac56-\frac45\right)^{-1}\cdot\left(\frac23\right)^{-1}+\left(\frac12\right)^{-1}$
  \task $\left(\frac52+\frac25\right)^{-1}:\left(\frac73\right)^{-1}-\left(\frac43\right)^2$
  \task $\left(\frac23+4\right):1,\widehat7-2$
  \task $\left(1-\frac34\right)^{-2}\cdot\left(3,\widehat{42}-1\right)+4$
  \task $\left(3-\frac45\cdot2,3\right)^{-1}+3,2\widehat4$
\end{tasks}\begin{tasks}[start=8](1)
  \task $\left\{\left[\left(-\frac35\right)^3\cdot\left(-\frac35\right)^2\right]^3:\left(-\frac35\right)^{15}\right\}-\left(\frac43\right)^3\cdot\left(\frac32\right)^4$
  \task $\left\{\left[\left(\frac13-1\right)\cdot\left(-1+\frac14\right)\right]^2:\left(2-\frac32\right)^2\right\}+\left[\left(-\frac15\right)^2\cdot\left(-\frac53\right)\right]\cdot\left(\frac54+10\right)$
  \task $\left[\left(\frac32-4\right):\left(1-\frac83\right)+\frac12\right]^4-\left(1+\frac23\right)^2\cdot\left(1-\frac25\right)^3$
\end{tasks}

\ex{1} Calcula las potencias de exponente negativo siguientes:\begin{tasks}(4)
  \task $3^{-2}$
  \task $(-3)^{-2}$
  \task $(-2)^{-1}$
  \task $(\sqrt3)^{-1}$
  \task $\left(-\frac12\right)^{-1}$
  \task $\left(\frac23\right)^{-2}$
  \task $\left(-\frac34\right)^{-2}$
  \task $\left(2-\frac12\right)^{-1}$
  \task $\left(\frac15-2\right)^{-2}$
  \task $\left(\frac{16}{5}-1,2\right)^{-3}$
  \task $\left(\frac34-\frac56\right)^{-2}$
  \task $\left(1,3-\frac45\right)^{-2}$
\end{tasks}

\ex{1} Calcula el valor de los radicales siguientes, justificando la respuesta:\begin{tasks}(5)
  \task $\sqrt{16}$
  \task $\sqrt[3]{-8}$
  \task $\sqrt[5]{32}$
  \task $\sqrt[3]{\frac{27}{8}}$
  \task $\sqrt[3]{0,064}$
\end{tasks}

\ex{1} Simplifica al máximo posible los radicales siguientes:\begin{tasks}(5)
  \task $\sqrt[3]{4a^4}$
  \task $\sqrt[15]{8a^3}$
  \task $\sqrt[6]{0,027}$
  \task $\sqrt[8]{16a^4}$
  \task $\sqrt[3]{a^2b^5}$
\end{tasks}

\ex{2} Reduce los radicales de cada apartado a mínimo común índice:\begin{tasks}(3)
  \task $\sqrt[3]{n^2},\ \sqrt[4]{m^3},\ \sqrt[6]{n^5}$
  \task $\sqrt{a^3},\ \sqrt[4]{m},\ \sqrt[5]{n^5},\ \sqrt[10]{a^3}$
  \task $\sqrt[3]{5xy^2},\ \sqrt{6x^3z},\ \sqrt[6]{\frac{3xy^3}{2z}}$
\end{tasks}

\ex{2} Realiza las operaciones indicadas:\begin{tasks}(2)
  \task $\sqrt{1+\sqrt{6+\sqrt{5+\sqrt{16}}}}$
  \task $\sqrt{3a^2+\sqrt{6a^4-\sqrt{25a^8}}}$
  \task $\sqrt{25\sqrt{81\sqrt{256}}}$
  \task $\left(\sqrt{a\sqrt{b\sqrt{c\sqrt d}}}\right)^{32}$
\end{tasks}

\ex{1} Calcula las raíces siguientes por el método más fácil:\begin{tasks}(3)
  \task $\sqrt[3]{8\cdot27\cdot64}$
  \task $\sqrt[5]{243\cdot32\cdot0,00001}$
  \task $\sqrt{49\cdot121\cdot169}$
  \task $\sqrt{25:0,0001}$
  \task $\sqrt[3]{8:0,064}$
  \task $\sqrt[4]{20736:256}$
  \task $\sqrt[5]{243:32}$
  \task $\sqrt{49^3}$
  \task $\sqrt[3]{0,064^2}$
  \task $\sqrt[3]{2197^4}$
  \task $\sqrt[4]{6561^3}$
  \task $\sqrt[5]{0,00032^4}$
\end{tasks}

\ex{2} Extrae todos los factores posibles de cada uno de los radicales siguientes:\begin{tasks}(4)
  \task $\sqrt8$
  \task $\sqrt{12}$
  \task $\sqrt[3]{16}$
  \task $\sqrt[3]{54}$
  \task $\sqrt[5]{64}$
  \task $\sqrt{\frac{27}{4}}$
  \task $\sqrt[5]{\frac{5x^{10}}{y^8}}$
  \task $\sqrt[3]{\frac{8x^4y^3z}{n^6}}$
  \task $3\sqrt{8a^3}$
  \task $2x^2y\sqrt{x^4y^3}$
  \task $\frac{xy^2}{3}\sqrt{27xy^3}$
  \task $\sqrt[4]{\frac{32x^6}{81y^5}}$
\end{tasks}

\ex{2} Introduce el coeficiente dentro del radical correspondiente y simplifica:\begin{tasks}(4)
  \task $7\sqrt a$
  \task $2a\sqrt{3a}$
  \task $3+\sqrt7$
  \task $x\sqrt{\frac1x}$
  \task $x^3y\sqrt{xy}$
  \task $\frac13\sqrt[4]{\frac{27}{2}}$
  \task $\frac32\sqrt{\frac23}$
  \task $\frac2a\sqrt{\frac{ax}{2}}$
  \task $\frac{3}{2xy}\sqrt{\frac{2xz}{y}}$
  \task $\frac23\sqrt[3]{\frac{81}{4}}$
  \task $\left(1+\frac12\right)\sqrt[3]{\frac{4}{81}}$
  \task $\frac{a-b}{a+b}\sqrt{\frac{a+b}{a-b}}$
  \task $(x+y)\sqrt{\frac{x-y}{x+y}}$
\end{tasks}

\ex{2} Racionaliza el denominador de las expresiones siguientes:\begin{tasks}(4)
  \task $\dfrac{5}{\sqrt5}$
  \task $\dfrac{2\sqrt6}{\sqrt2}$
  \task $\sqrt{\dfrac52}$
  \task $\dfrac{3}{\sqrt5}$
  \task $\sqrt[3]{\dfrac32}$
  \task $\dfrac{3}{2\sqrt[3]4}$
  \task $\dfrac{3y}{x\sqrt y}$
  \task $\dfrac{x\sqrt y}{y\sqrt x}$
  \task $\dfrac{2\sqrt{12}}{5\sqrt3}$
  \task $\dfrac{3\sqrt2}{\sqrt6}$
  \task $\dfrac{3xy^2}{\sqrt[3]{x^2y}}$
  \task $\dfrac{6(3-y)}{\sqrt[3]{(3-y)^2}}$
  \task $\dfrac{-4}{\sqrt[4]3\cdot\sqrt[3]2}$
\end{tasks}

\ex{2} Racionaliza el denominador de las expresiones siguientes:\begin{tasks}(4)
  \task $\dfrac{1}{\sqrt2+\sqrt3}$
  \task $\dfrac{2}{\sqrt2-1}$
  \task $\dfrac{\sqrt2}{3-\sqrt2}$
  \task $\dfrac{1-\sqrt2}{1+\sqrt2}$
  \task $\dfrac{\sqrt5-2}{3-2\sqrt5}$
  \task $\dfrac{1}{\sqrt x-\sqrt y}$
  \task $\dfrac{\sqrt{2y}}{\sqrt2-\sqrt y}$
  \task $\dfrac{\sqrt{xy}}{2\sqrt x-3\sqrt y}$
  \task $\dfrac{5}{2(\sqrt3-\sqrt2)}$
  \task $\dfrac{\sqrt2}{2\sqrt2+1}$
  \task $\dfrac{2+\sqrt2}{2-\sqrt2}$
  \task $\dfrac{-\sqrt2}{\sqrt[3]2(\sqrt{125}+2)}$
  \task $\dfrac{1}{\sqrt{3+\sqrt6}}$
\end{tasks}\recorda{Si en el denominador hay una suma o una resta con raíces, multiplicamos por el conjugado: $(a+b)(a-b)=a^2-b^2$.}

\ex{2} Comprueba que los radicales de cada apartado son semejantes:\begin{tasks}(2)
  \task $5\sqrt2,\ 3\sqrt8,\ \sqrt{18}$
  \task $\sqrt{3a^2},\ \sqrt{27},\ \sqrt{\frac13},\ 2\sqrt{\frac{4}{27}}$
  \task $\sqrt[3]{81},\ \sqrt[3]{24},\ \sqrt[6]9,\ \sqrt[3]{\frac19a^3}$
  \task $3\sqrt[4]{2xy},\ \sqrt[4]{32x^5y},\ \sqrt[4]{2x^9y^5}$
\end{tasks}

\ex{1} Calcula y simplifica:\begin{tasks}(3)
  \task $\sqrt{25:0,01}$
  \task $\sqrt[4]{9x^5y^3}:\sqrt[4]{3x^3y}$
  \task $\sqrt[3]{(27x^3y^6)^5}$
\end{tasks}

\ex{2} Realiza las sumas y restas siguientes:\begin{tasks}(2)
  \task $\sqrt8+\sqrt{18}+\sqrt{32}+\sqrt{50}$
  \task $\sqrt{28}+\sqrt{63}-\sqrt{175}$
  \task $6\sqrt3-4\sqrt3+5\sqrt3$
  \task $3\sqrt2+5\sqrt2-8\sqrt2$
  \task $3\sqrt2-3\sqrt8+3\sqrt{18}$
  \task $2a\sqrt3-\sqrt{27a^2}+a\sqrt{12}$
  \task $(3+a)\sqrt5-\sqrt{125}+\sqrt{5a^2}$
  \task $7\sqrt{54}-3\sqrt{18}+\sqrt{24}-\frac35\sqrt{50}-\sqrt6$
  \task $\sqrt{\frac13}+\sqrt{27}$
  \task $\sqrt{\frac12}+\sqrt2+\sqrt[4]4+\sqrt[6]8+\sqrt[4]{64}$
  \task $2\sqrt{\frac34}-\sqrt{27}+\frac14\sqrt{12}-3\sqrt{\frac{75}{9}}$
  \task $\sqrt{18}+5\sqrt{12}-2\sqrt2-\sqrt{243}$
  \task $2\sqrt3+5\sqrt[6]{27}-3\sqrt{12}$
  \task $5\sqrt{20}-2\sqrt{45}+\sqrt5$
\end{tasks}

\ex{2} Realiza las operaciones siguientes:\begin{tasks}(2)
  \task $\sqrt5\cdot\sqrt6$
  \task $\sqrt5\cdot\sqrt[3]6$
  \task $(\sqrt2)^2+(\sqrt2)^3+(\sqrt2)^4$
  \task $a\cdot\sqrt a\cdot\sqrt[3]a\cdot\sqrt[4]a$
  \task $(\sqrt3)^5-(\sqrt3)^3-(2\sqrt3)^2$
  \task $2\sqrt3\left(\sqrt[3]4-\sqrt[3]3\right)$
  \task $\sqrt{\frac23}\cdot\sqrt[3]{\frac23}\cdot\sqrt[4]{\frac32}$
  \task $\sqrt{12}\cdot\sqrt{\frac34}\cdot\sqrt{\frac{12}{5}}\cdot\sqrt{\frac{15}{4}}$
  \task $\sqrt{10}\cdot\sqrt[3]{10}\cdot\sqrt[4]{10}$
  \task $(\sqrt3+\sqrt2)^2$
\end{tasks}

\ex{2} Escribe las expresiones siguientes bajo un solo radical y simplifica los resultados:\begin{tasks}(3)
  \task $\sqrt{\sqrt8}$
  \task $\sqrt{2\sqrt2}$
  \task $\sqrt[3]{\sqrt{32}}$
  \task $\sqrt{2\sqrt{3\sqrt[4]4}}$
  \task $\sqrt[3]{5\sqrt[4]{5\sqrt[5]{5^2}}}$
  \task $\sqrt[3]{3\sqrt[3]{\frac19}}$
\end{tasks}


% =====================================================
\seccio{5}{Logaritmos}

\exf{1} Aplicando la definición, calcula el valor de los logaritmos siguientes:\begin{tasks}(4)
  \task $\log_2 8$
  \task $\log_3 81$
  \task $\log 1000$
  \task $\log 0,0001$
  \task $\ln e^{33}$
  \task $\ln e^{-4}$
  \task $\log_4 16$
  \task $\log_4 0,25$
\end{tasks}

\ex{1} Aplicando la definición, calcula el valor de los logaritmos siguientes:\begin{tasks}(4)
  \task $\log_3\sqrt{27}$
  \task $\log_7\frac{1}{49}$
  \task $\log_3 3\sqrt3$
  \task $\log_3 0,\widehat3$
  \task $\log_5\sqrt[3]{25}$
  \task $\log_9\sqrt[3]3$
  \task $\log_{\frac12}\sqrt8$
  \task $\log_8 0,125$
\end{tasks}

\ex{2} Calcula, si es posible, el valor de $x$ en cada una de las expresiones siguientes:\begin{tasks}(4)
  \task $\log_x 8=-3$
  \task $\log_{-3}\sqrt x=9$
  \task $\log_3(-81)=x$
  \task $\log_{\frac{1}{\sqrt2}}x=-2$
  \task $\log_x\sqrt2=0$
  \task $\log_1 2=x$
  \task $\log_5 x=-1$
  \task $\log_{\frac1a}a^2=x$
\end{tasks}\recorda{$\log_a b$ solo existe si $a>0$, $a\neq1$ y $b>0$.}

\ex{1} Halla el resultado de las expresiones siguientes mediante las propiedades de los logaritmos:\begin{tasks}(1)
  \task $2\log_4 16+\log_2 32-3\log_7 49$
  \task $\log_2 8+\log_3 27+3\log_5 125$
  \task $\log_5 625-\log_9 81+\log_8 64$
\end{tasks}

\ex{2} Una bacteria se duplica cada hora si tiene las condiciones óptimas. ¿Cuántas horas tendrán que pasar para que una bacteria inicial cree una colonia de mil millones de bacterias?

\ex{2} Considerando $\log2=0,301$, $\log3=0,477$ y $\log k=0,778$, calcula los logaritmos siguientes:\begin{tasks}(3)
  \task $\log50$
  \task $\log0,\widehat3$
  \task $\log_k3$
  \task $\log\sqrt{12k}$
  \task $\log_2\sqrt[3]{6k}$
  \task $\log4,5$
\end{tasks}

\ex{2} Si $\log_3 2=0,631$ y $\log_3 5=1,465$, calcula el valor del logaritmo en base 3 de 150.

\ex{2} Toma logaritmos decimales en las igualdades siguientes y desarrolla las expresiones:\begin{tasks}(3)
  \task $P=10x^3yz^3$
  \task $Q=\dfrac{100x^2}{x+y}$
  \task $R=\sqrt[3]{\dfrac{2x^2y^5}{3z^3}}$
  \task $x=a^4b^3c^{\frac32}$
  \task $y=\dfrac{\sqrt[3]{x^2}}{ax}$
  \task $xy=\dfrac{(m+2n)n^2}{m-2n}$
\end{tasks}

\ex{2} Expresa el valor de $E$ en cada caso sin que aparezcan logaritmos:\begin{tasks}(1)
  \task $\log E=2-2\log x+\log y-5\log z$
  \task $\log E=3\log2-4\log x+3\log y-2\log z$
  \task $\log E=\log(x-2y)+\log(x+y)$
  \task $\log E=3\log(x+10)-\log\frac{2x+20}{3}+\log\frac23$
\end{tasks}

\ex{1} Con la ayuda de la calculadora, obtén aproximaciones hasta las milésimas de los logaritmos siguientes:\begin{tasks}(3)
  \task $\log_3 20$
  \task $\log_{\frac14}\frac75$
  \task $\log_{0,5}60$
  \task $\log_{\sqrt2}3$
  \task $\log_{\frac25}\sqrt[3]2$
  \task $\log_{\sqrt2}\sqrt3$
\end{tasks}\recorda{Cambio de base: $\log_a x=\dfrac{\log x}{\log a}=\dfrac{\ln x}{\ln a}$.}

\ex{2} Determina el valor de $x$:\begin{tasks}(4)
  \task $8^x=1024$
  \task $3^{x^2}=27$
  \task $3^{x^2-6}=27$
  \task $10^{x-1}=10^3$
  \task $8^{x-3}=1024$
  \task $(3^x)^2=27$
  \task $3^{x^2}+18=27$
  \task $2^{x^2-2x+1}=1$
\end{tasks}


% =====================================================
\seccio{6}{Ampliación y problemas}

\exf{2} Clasifica cada número en el menor conjunto numérico. ¡Opera primero!\begin{tasks}(3)
  \task $\sqrt2\cdot\sqrt8$
  \task $\sqrt2+\sqrt8$
  \task $\pi-\pi$
  \task $0,1010010001\ldots$
  \task $-\sqrt{\frac94}$
  \task $\sqrt{-4}$
\end{tasks}

\ex{1} Halla la fracción generatriz y simplifícala:\begin{tasks}(4)
  \task $0,1\widehat6$
  \task $2,\widehat{09}$
  \task $-1,2\widehat{45}$
  \task $0,\widehat9$
\end{tasks}

\ex{2} Expresa cada conjunto con una desigualdad con valor absoluto (busca el centro y el radio):\begin{tasks}(4)
  \task $(1,\, 7)$
  \task $[-5,\, 1]$
  \task $(-\infty,\, -1)\cup(5,\, +\infty)$
  \task $[-0,5;\ 0,5]$
\end{tasks}

\ex{2} La temperatura $T$ de una nevera debe ser de $4\,^\circ$C con un margen de error de $1,5\,^\circ$C. Escribe esta condición con un valor absoluto y como intervalo. ¿Es aceptable una lectura de $5,7\,^\circ$C?

\ex{1} Entornos:\begin{tasks}(2)
  \task Escribe como intervalo $B(-1;\ 2,5)$.
  \task Escribe como intervalo el entorno cerrado de centro 3 y radio $\frac12$.
  \task ¿Centro y radio del intervalo $(-2,\, 8)$?
  \task ¿Qué enteros pertenecen a $B(0,5;\ 3)$?
\end{tasks}

\ex{2} Un cuadrado tiene un área de $50\ \text{cm}^2$. Calcula su lado y su diagonal en forma de radical simplificado.

\ex{2} Racionaliza, opera y simplifica:\begin{tasks}(2)
  \task $\dfrac{\sqrt3+1}{\sqrt3-1}-\dfrac{\sqrt3-1}{\sqrt3+1}$
  \task $\dfrac{1}{\sqrt2}+\dfrac{1}{\sqrt2+1}$
\end{tasks}

\ex{2} Calcula sin calculadora pasando a exponentes fraccionarios:\begin{tasks}(3)
  \task $\dfrac{8^{2/3}\cdot27^{-1/3}}{16^{3/4}}$
  \task $\left(\dfrac{1}{32}\right)^{-2/5}$
  \task $81^{0,75}$
\end{tasks}

\ex{1} Encuentra el error en cada uno de estos cálculos y corrígelo:\begin{tasks}(2)
  \task $\sqrt{9+16}=3+4=7$
  \task $-3^2=9$
  \task $\log(5+5)=\log5+\log5$
  \task $2^{-3}=-8$
\end{tasks}

\ex{2} El pH de una disolución es $\text{pH}=-\log[\text{H}^+]$.\begin{tasks}(1)
  \task Calcula el pH si $[\text{H}^+]=3,2\cdot10^{-4}$ mol/L.
  \task La sangre tiene un pH de 7,4. ¿Cuál es su concentración $[\text{H}^+]$?
\end{tasks}

\ex{2} Ingresamos 2000 € en una cuenta al 3\,\% de interés compuesto anual. ¿Cuántos años tienen que pasar, como mínimo, para que haya 3000 €?

\ex{3} Sin calculadora, ordena de menor a mayor $\log_2 10$, \ $3$ \ y \ $\log_3 10$. Justifícalo.

\vfill
\begin{center}\color{muted}\small\faCheckCircle[regular]\ \ ¡Ánimo! Comprueba tus resultados con las herramientas interactivas de los apuntes de la unidad.\end{center}

\end{document}
